%% kapa: exam=1 round=31 year=2020 period=1 form=A subject=civil_law subject_ko=민법 seq=12 number=12
%% source: Q-Net 감정평가사 기출문제 (cst00309) artlSeq=5207219 file=extracted/1008/5207219/A형/A형/2020년도 제31회 감정평가사 1차 시험_1교시_A형.pdf
%% conversion: kordoc 4.24.51 -> tools/gen1.py (auto); review=proofread-sonnet; flags=none
\begin{question}{12}
착오에 관한 설명으로 옳지 않은 것은? (다툼이 있으면 판례에 따름)
\begin{choices}
  \item 매도인이 매매대금 미지급을 이유로 매매계약을 해제한 후에도 매수인은 착오를 이유로 이를 취소할 수 있다.
  \item 보험회사의 설명의무 위반으로 보험계약의 중요사항을 제대로 이해하지 못하고 착오에 빠져 계약을 체결한 고객은 그 계약을 취소할 수 있다.
  \item 계약서에 X토지를 목적물로 기재한 때에도 Y토지에 대하여 의사의 합치가 있었다면 Y토지를 목적으로 하는 계약이 성립한다.
  \item 착오에 관한 민법규정은 법률의 착오에 적용되지 않는다.
  \item 취소의 의사표시는 취소자가 그 착오를 이유로 자신의 법률행위의 효력을 처음부터 없애려는 의사가 드러나면 충분하다.
\end{choices}
\end{question}
\answer{4}
